\section{Discussion and Future Work}
\label{sec:disc}

\subsection{Answer to Research Questions}~\label{sec:disc:rq}

We answer the research questions as follows:

\begin{itemize}
  \item \textbf{RQ1:} Our parallel congruence closure algorithm achieves a
  sublinear, but meaningful speedup on random benchmarks, a synthetic benchmark
  suite designed to stress test our algorithm, and a set of equality saturation
  benchmarks~\cite{eqsat}. We observe meaningful speedups up to 32 cores, but we
  experienc a pleateau on 32 and then 64 cores.
  \item \textbf{RQ2:} Our algorithm tends to perform better on benchmarks with a
  large number of merges per round. Having more rounds of congruence closure
  does not necessarily lead to worse performance, but it does not lead to better
  performance either as we only parallelize within rounds.
  \item \textbf{RQ3:} Our algorithm provides a meaningful speedup over a
  sequential implementation on the egg benchmarks. Future work should evaluate
  our algorithm on benchmarks from SMT solvers and  more challenging equality
  saturation benchmarks such as \texttt{egg-cc}~\cite{eggcc}.
\end{itemize}


\subsection{Limitations of Evaluation}~\label{sec:disc:limit-eval}

There are several limitations of our evaluation that we hope to address in
future work.

\textbf{Machine Setup.} We had hoped to run our final benchmarks on a modern
Amazon EC2 \texttt{c8i.metal-48xl} instance (on which we had obtained more
promising preliminary results), but were limited by time and cost to run our
final benchmark suite on this machine. Our final benchmark suite took over 6
hours to complete, which would have cost about \$60 on the EC2 instance.

\textbf{Workloads.} We hope in the future to evaluate our implementation on a
wider-reaching and more realistic set of benchmarks. For example, we hope
evaluate our implementation on SMT-LIB benchmarks, which requires support for
disjunctions in forumulas, e.g. those of the form $x = y \vee w = z$.

In order to accomodate disjunctions, we would have to either implement on our
algorithm in a modern SMT solver, which is difficult because of limitations we
discuss in \autoref{sec:disc:limit-alg}. Alternativley, we could run an existing
SMT solver on the SMT-LIB benchmarks to obtain a set of equalities and
disequalities that are relevant to the satisfiability of the formula, and then
run our algorithm on these equalities and disequalities. This is future work.

\textbf{Comparisons to Baseline.} We have implemented a very simple sequential
baseline based on the classic Nelson-Oppen algorithm. In the future, we hope to
extend our evaluation to compare to state-of-the-art sequential solvers, like
those in \texttt{egg} and \texttt{egglog}. 
% It is not quite an apples-to-apples comparison, since these solvers support
% more features (like dynamic addition of e-nodes) that we do not.

\subsection{Limitations of Algorithm and Future Work}~\label{sec:disc:limit-alg}

There are several limitations to the algorithm as presented that we would like
to address in future work. 

\textbf{Synchronous Round Structure.} Our synchronous round structure inherently
limits parallel scaling. In particular, the instance size processed at each
round is small relative to the size of the problem for several benchmarks we
have examined (like the \texttt{egg} benchmarks). We will likely have to take a
more asynchronous approach to achieve bett thread scaling; this would entail
generating new congruences from discovered equalities immediately instead of
waiting for the round to complete.

\textbf{Backtracking.} Our approach is not \textit{backtrackable}, meaning that
it is not possible to derive an extractable proof how why a congruence holds
(i.e. the proof tree). Implementing backtrackability requires maintaining the
path along the union find that represents a chain of equalities, which is
extremely difficult in a concurrent setting when updates arrive asynchronously
and the path is constantly changing. It may be possible to use snapshotting
techniques in order to derive this path.

\textbf{Proof Production.} Our implementation does not produce \emph{proofs} or
\emph{witnesses} of the equalities it discovers. For instance if our algorithm
discovers that $x = y$, it cannot replay the chain of equalities and congruences
that led to this conclusion. This is important for a number of reasons: it is
necessary to justify correctness and it is used to produce conflict clauses in
SMT solvers.

Finally, we would like to support our implementation to incremental updates
(where equalities and new E-node may arrive dynamically). This is important for
both equality saturation and SMT solvers, where new equalities may be discovered
via for example quantifier instantiation.
